\chapter{2012年浙江卷（文）}

\section{单选题}

\begin{problem}
设全集 \(U = \{1, 2, 3, 4, 5, 6\}\)，集合 \(P = \{1, 2, 3, 4\}\)，\(Q = \{3, 4, 5\}\)，则 \(P \cap (\complement_U Q) =\)（\quad）

\choices
    {\(\{1, 2, 3, 4, 6\}\)}
    {\(\{1,2,3,4,5\}\)}
    {\(\{1,2,5\}\)}
    {\(\{1,2\}\)}
\end{problem}

\begin{answer}
D
\end{answer}

\begin{solution}
全集中不属于 \(Q\) 的元素为
\[
\complement_U Q=\{1,2,6\}
\]
因此
\[
P\cap(\complement_U Q)=\{1,2,3,4\}\cap\{1,2,6\}=\{1,2\}
\]
故选 D.
\end{solution}


\begin{problem}
已知 \(\i\) 是虚数单位，则 \( \frac{3+\i}{1-\i}= \)（\quad）

\choices
    {\(1 - 2\i\)}
    {\(2 - \i\)}
    {\(2 + \i\)}
    {\(1 + 2\i\)}
\end{problem}

\begin{answer}
D
\end{answer}

\begin{solution}
分子分母同乘 \(1+\i\)，得
\[
\frac{3+\i}{1-\i}
=\frac{(3+\i)(1+\i)}{(1-\i)(1+\i)}
=\frac{2+4\i}{2}=1+2\i
\]
故选 D.
\end{solution}

\begin{problem}
已知某三棱锥的三视图（单位：\(\mathrm{cm}\)）如图所示，则该三棱锥的体积是（\quad）

\bitmapfigure[max width=0.42\linewidth]{img/2012/zhejiang_liberal/q3_fig1.png}

\choices
    {\(1 \mathrm{~cm}^{3}\)}
    {\(2 \mathrm{~cm}^{3}\)}
    {\(3 \mathrm{~cm}^{3}\)}
    {\(6 \mathrm{~cm}^{3}\)}
\end{problem}

\begin{answer}
A
\end{answer}

\begin{solution}
由三视图可知，俯视图给出底面是直角三角形，两个直角边长分别为 \(1\) 和 \(2\)；正视图给出三棱锥的高为 \(3\). 所以
\[
V=\frac{1}{3}\times\frac{1}{2}\times1\times2\times3=1\ \mathrm{cm}^{3}
\]
故选 A.
\end{solution}

\begin{problem}
设 \(l\) 是直线，\(\alpha\)，\(\beta\) 是两个不同的平面，下列选项正确的是（\quad）

\choices
    {若 \(l \parallel \alpha\)，\(l \parallel \beta\)，则 \(\alpha \parallel \beta\)}
    {若 \(l \parallel \alpha\)，\(l \perp \beta\)，则 \(\alpha \perp \beta\)}
    {若 \(\alpha \perp \beta\)，\(l \perp \alpha\)，则 \(l \perp \beta\)}
    {若 \(\alpha \perp \beta\)，\(l \parallel \alpha\)，则 \(l \perp \beta\)}
\end{problem}

\begin{answer}
B
\end{answer}

\begin{solution}
选项 B 正确. 因为 \(l\parallel\alpha\)，过 \(\alpha\) 内任一点作直线 \(m\parallel l\)，则 \(m\subset\alpha\). 又 \(l\perp\beta\)，所以 \(m\perp\beta\)，从而平面 \(\alpha\) 内有一条直线垂直于平面 \(\beta\)，故 \(\alpha\perp\beta\).

其余选项均不成立. 例如取 \(\alpha:z=0\)、\(\beta:y=0\)：取方向为 \(x\) 轴方向且位于 \(y=z=1\) 的直线，可同时平行于 \(\alpha\)，\(\beta\)，但两平面不平行，故 A 错；取方向为 \(z\) 轴方向且位于 \(x=y=1\) 的直线，虽有 \(l\perp\alpha\) 且 \(\alpha\perp\beta\)，但 \(l\parallel\beta\)，故 C 错；取方向为 \(x\) 轴方向且位于 \(y=0\)，\(z=1\) 的直线，虽有 \(l\parallel\alpha\) 且 \(\alpha\perp\beta\)，但 \(l\) 不垂直于 \(\beta\)，故 D 错.
\end{solution}

\begin{problem}
设 \(a \in \R\)，则 \(a = 1\) 是“直线 \(l_1\)：\(ax + 2y - 1 = 0\) 与直线 \(l_2\)：\(x + 2y + 4 = 0\) 平行”的（\quad）

\choices
    {充分不必要条件}
    {必要不充分条件}
    {充分必要条件}
    {既不充分也不必要条件}
\end{problem}

\begin{answer}
C
\end{answer}

\begin{solution}
两直线的法向量分别为 \((a,2)\) 和 \((1,2)\). 两直线平行的必要条件是法向量平行，因此
\[
\begin{vmatrix}a&2\\1&2\end{vmatrix}=2a-2=0
\]
从而 \(a=1\). 此时两直线分别为 \(x+2y-1=0\) 与 \(x+2y+4=0\)，法向量相同而常数项不同，故两直线确实平行. 因此 \(a=1\) 是该两直线平行的充分必要条件，故选 C.
\end{solution}

\begin{problem}
把函数 \(y = \cos 2x + 1\) 的图象上所有点的横坐标伸长到原来的 \(2\) 倍（纵坐标不变），然后向左平移 \(1\) 个单位长度，再向下平移 \(1\) 个单位长度，得到的图象是（\quad）

\choices
    {\choicebitmap[width=0.15\paperwidth]{img/2012/zhejiang_liberal/q6_A.png}}
    {\choicebitmap[width=0.15\paperwidth]{img/2012/zhejiang_liberal/q6_B.png}}
    {\choicebitmap[width=0.15\paperwidth]{img/2012/zhejiang_liberal/q6_C.png}}
    {\choicebitmap[width=0.15\paperwidth]{img/2012/zhejiang_liberal/q6_D.png}}
\end{problem}

\begin{answer}
A
\end{answer}

\begin{solution}
设原图象上的点为 \((u,\cos2u+1)\). 横坐标伸长为原来的 \(2\) 倍后变为 \((2u,\cos2u+1)\)，再向左平移 \(1\) 个单位、向下平移 \(1\) 个单位，得到
\[
(x,y)=(2u-1,\cos2u)
\]
由 \(u=(x+1)/2\)，新图象方程为
\[
y=\cos(x+1)
\]
它的最大值为 \(1\)，在 \(x=-1\) 处取得；在 \(x=\frac{\pi}{2}-1\) 处过 \(x\) 轴. 与选项图象相符的是 A.
\end{solution}

\begin{problem}
设 \(\bs{a}\)，\(\bs{b}\) 是两个非零向量（\quad）

\choices
    {若 \(|\bs{a} + \bs{b}| = |\bs{a}| - |\bs{b}|\)，则 \(\bs{a} \perp \bs{b}\)}
    {若 \(\bs{a} \perp \bs{b}\)，则 \(|\bs{a} + \bs{b}| = |\bs{a}| - |\bs{b}|\)}
    {若 \(|\bs{a} + \bs{b}| = |\bs{a}| - |\bs{b}|\)，则存在实数 \(\lambda\)，使得 \(\bs{b} = \lambda \bs{a}\)}
    {若存在实数 \(\lambda\)，使得 \(\bs{b} = \lambda \bs{a}\)，则 \(|\bs{a} + \bs{b}| = |\bs{a}| - |\bs{b}|\)}
\end{problem}

\begin{answer}
C
\end{answer}

\begin{solution}
由已知等式两边平方，得
\[
|\bs{a}|^{2}+|\bs{b}|^{2}+2\bs{a}\cdot\bs{b}
=(|\bs{a}|-|\bs{b}|)^{2}
=|\bs{a}|^{2}+|\bs{b}|^{2}-2|\bs{a}||\bs{b}|
\]
所以
\[
\bs{a}\cdot\bs{b}=-|\bs{a}||\bs{b}|
\]
由反向三角不等式的等号条件，\(\bs{a}\) 与 \(\bs{b}\) 反向共线，故存在实数 \(\lambda\) 使 \(\bs{b}=\lambda\bs{a}\)，所以 C 正确. A 不成立，例如 \(\bs{b}=-\bs{a}\) 时已知等式成立但两向量不垂直；B 不成立，例如两向量垂直且等长时右边为 \(0\) 而左边为 \(\sqrt2|\bs{a}|\)；D 不成立，例如 \(\bs{b}=2\bs{a}\) 时等式两边分别为 \(3|\bs{a}|\) 和 \(-|\bs{a}|\).
\end{solution}

\begin{problem}
如图，中心均为原点 \(O\) 的双曲线与椭圆有公共焦点，\(M\)、\(N\) 是双曲线的两顶点. 若 \(M\)、\(O\)、\(N\) 将椭圆长轴四等分，则双曲线与椭圆的离心率的比值是（\quad）

\FigureLayoutDeclare{img/2012/zhejiang_liberal/q8_fig1.png}{7.4cm}{18bp 50bp 113bp 49bp}
\bitmapfigure[max width=0.52\linewidth]{img/2012/zhejiang_liberal/q8_fig1.png}

\choices
    {\(3\)}
    {\(2\)}
    {\(\sqrt{3}\)}
    {\(\sqrt{2}\)}
\end{problem}

\begin{answer}
B
\end{answer}

\begin{solution}
设椭圆长半轴为 \(A\)，焦半距为 \(c\)，双曲线实半轴为 \(a\). 由于双曲线两顶点将椭圆长轴四等分，双曲线顶点到原点的距离为 \(A/2\)，故
\[
a=\frac A2
\]
两曲线公共焦点，故椭圆离心率为 \(e_1=\frac cA\)，双曲线离心率为 \(e_2=\frac ca\).
于是
\[
\frac{e_2}{e_1}=\frac{c/a}{c/A}=\frac Aa=2
\]
故选 B.
\end{solution}

\begin{problem}
若正数 \(x\)，\(y\) 满足 \(x + 3y = 5xy\)，则 \(3x + 4y\) 的最小值是（\quad）

\choices
    {\( \frac{24}{5} \)}
    {\( \frac{28}{5} \)}
    {\(5\)}
    {\(6\)}
\end{problem}

\begin{answer}
C
\end{answer}

\begin{solution}
由 \(x+3y=5xy\) 且 \(x\)，\(y>0\)，可得
\[
\frac3x+\frac1y=5
\]
令 \(u=3/x>0\)，\(v=1/y>0\)，则 \(u+v=5\)，而
\[
3x+4y=\frac9u+\frac4v
\]
由柯西不等式，
\[
\left(\frac9u+\frac4v\right)(u+v)\geq(3+2)^2=25
\]
因为 \(u+v=5\)，所以 \(3x+4y\geq5\). 当且仅当 \(3/u=2/v\)，即 \(u:v=3:2\) 时取等，此时 \(x=1\)，\(y=\frac12\)，满足原条件. 故最小值为 \(5\)，选 C.
\end{solution}

\begin{problem}
设 \(a > 0\)，\(b > 0\)，\(\e\) 是自然对数的底数（\quad）

\choices
    {若 \(\e^a + 2a = \e^b + 3b\)，则 \(a > b\)}
    {若 \(\e^a + 2a = \e^b + 3b\)，则 \(a < b\)}
    {若 \(\e^a - 2a = \e^b - 3b\)，则 \(a > b\)}
    {若 \(\e^a - 2a = \e^b - 3b\)，则 \(a < b\)}
\end{problem}

\begin{answer}
A
\end{answer}

\begin{solution}
若 \(a\leq b\)，则 \(\e^a\leq\e^b\)，从而
\[
\e^a+2a\leq\e^b+2b<\e^b+3b
\]
这与题设等式矛盾. 因此必有 \(a>b\)，故选 A.
\end{solution}

\section{填空题}

\begin{problem}
某个年级有男生 \(560\text{ 人}\)，女生 \(420\text{ 人}\)，用分层抽样的方法从该年级全体学生中抽取一个容量为 \(280\) 的样本，则此样本中男生人数为\fillinblank{}.
\end{problem}

\begin{answer}
\(160\)
\end{answer}

\begin{solution}
全体学生人数为 \(560+420=980\)，男生所占比例为 \(\frac{560}{980}=\frac47\). 按分层抽样的比例抽取，样本中男生人数为
\[
280\times\frac{560}{980}=280\times\frac47=160
\]
\end{solution}

\begin{problem}
从边长为 \(1\) 的正方形的中心和顶点这五点中，随机（等可能）取两点，则该两点间的距离为 \(\frac{\sqrt{2}}{2}\) 的概率是\fillinblank{}.
\end{problem}

\begin{answer}
\(\frac{2}{5}\)
\end{answer}

\begin{solution}
五点中任取两点共有 \(\mathrm C_5^2=10\) 种等可能结果. 中心到四个顶点的距离均为正方形对角线的一半，即 \(\sqrt2/2\)，共 \(4\) 种；任意两个顶点间的距离为 \(1\) 或 \(\sqrt2\)，不符合条件. 因此所求概率为
\[
\frac4{10}=\frac25
\]
\end{solution}

\begin{problem}
若某程序框图如图所示，则该程序运行后输出的值是\fillinblank{}.

\texfigure[max width=0.42\linewidth]{img_repaint/2012/zhejiang_liberal/q13_fig1.tex}
\end{problem}

\begin{answer}
\(\frac{1}{120}\)
\end{answer}

\begin{solution}
程序先赋值 \(T=1\)，\(i=1\)，然后重复执行 \(T\leftarrow T/i\)、\(i\leftarrow i+1\)，直到 \(i>5\) 时输出 \(T\). 因此依次除以 \(1\)，\(2\)，\(3\)，\(4\)，\(5\)，输出值为
\[
T=1\div1\div2\div3\div4\div5=\frac1{120}
\]
\end{solution}

\begin{problem}
设 \(z = x + 2y\)，其中实数 \(x\)，\(y\) 满足 \(\begin{cases}
x - y + 1 \geqslant 0,\\
x + y - 2 \leqslant 0,\\
x \geqslant 0,\\
y \geqslant 0,
\end{cases}\) 则 \(z\) 的取值范围是\fillinblank{}.
\end{problem}

\begin{answer}
\(\left[0, \frac{7}{2}\right]\)
\end{answer}

\begin{solution}
约束条件等价于
\[
\begin{cases}
x\geq0,\\
y\geq0,\\
y\leq x+1,\\
y\leq2-x.
\end{cases}
\]
可行域的顶点为 \((0,0)\)，\((2,0)\)，\((0,1)\)，\(\left(\frac12,\frac32\right)\)，其中最后一个点由 \(x+1=2-x\) 求得. 在线性目标函数 \(z=x+2y\) 的可行域上，最大值和最小值均在顶点处取得. 逐点计算：\(z(0,0)=0\)，\(z(2,0)=2\)，\(z(0,1)=2\)，\(z\left(\frac12,\frac32\right)=\frac12+3=\frac72\).
故
\[
z\in\left[0,\frac72\right]
\]
\end{solution}

\begin{problem}
在 \(\triangle ABC\) 中，\(M\) 是 \(BC\) 的中点，\(AM = 3\)，\(BC = 10\)，则 \(\overrightarrow{AB} \cdot \overrightarrow{AC} = \)\fillinblank{}.
\end{problem}

\begin{answer}
\(-16\)
\end{answer}

\begin{solution}
由 \(M\) 是 \(BC\) 的中点，得
\[
\overrightarrow{AM}=\frac{\overrightarrow{AB}+\overrightarrow{AC}}2
\]
所以
\[
|\overrightarrow{AB}+\overrightarrow{AC}|=2AM=6
\]
又
\[
|\overrightarrow{AC}-\overrightarrow{AB}|=|BC|=10
\]
分别平方并相减：
\[
|\overrightarrow{AC}-\overrightarrow{AB}|^2-
|\overrightarrow{AB}+\overrightarrow{AC}|^2
=-4\overrightarrow{AB}\cdot\overrightarrow{AC}=100-36=64
\]
因此
\[
\overrightarrow{AB}\cdot\overrightarrow{AC}=-16
\]
\end{solution}

\begin{problem}
设函数 \( f(x) \) 是定义在 \( \R \) 上的周期为 2 的偶函数，当 \( x \in [0,1] \) 时，\( f(x) = x + 1 \)，则 \( f\left(\frac{3}{2}\right) = \)\fillinblank{}.
\end{problem}

\begin{answer}
\(\frac{3}{2}\)
\end{answer}

\begin{solution}
由周期性，\(f(3/2)=f(-1/2)\)，因为周期为 \(2\). 又因 \(f\) 为偶函数，
\[
f\left(-\frac12\right)=f\left(\frac12\right)
\]
由于 \(1/2\in[0,1]\)，所以
\[
f\left(\frac12\right)=\frac12+1=\frac32
\]
故 \(f(3/2)=3/2\).
\end{solution}

\begin{problem}
定义：曲线 \(C\) 上的点到直线 \(l\) 的距离的最小值称为曲线 \(C\) 到直线 \(l\) 的距离. 已知曲线 \(C_1\)：\(y = x^{2} + a\) 到直线 \(l\)：\(y = x\) 的距离等于曲线 \(C_2\)：\(x^{2} + (y + 4)^{2} = 2\) 到直线 \(l\)：\(y = x\) 的距离，则实数 \(a =\fillinblank{}\).
\end{problem}

\begin{answer}
\(\frac{9}{4}\)
\end{answer}

\begin{solution}
直线 \(y=x\) 写成 \(x-y=0\). 曲线 \(C_2\) 是圆心 \(C(0,-4)\)、半径 \(r=\sqrt2\) 的圆. 圆心到直线的距离为
\[
\frac{|0-(-4)|}{\sqrt2}=2\sqrt2
\]
故圆到直线的最小距离为
\[
2\sqrt2-\sqrt2=\sqrt2
\]

对 \(C_1\) 上点 \((x,x^2+a)\)，其到直线的距离为
\[
\frac{|x-(x^2+a)|}{\sqrt2}
=\frac{|-\left(x-\frac12\right)^2+\frac14-a|}{\sqrt2}
\]
若 \(a\leq1/4\)，括号内的表达式能取到 \(0\)，曲线与直线相交，距离为 \(0\)，不符合已知条件. 因此 \(a>1/4\)，此时括号内恒为负，故
\[
d(C_1,l)=\frac{a-\frac14}{\sqrt2}
\]
令其等于 \(\sqrt2\)，得
\(\frac{a-\frac14}{\sqrt2}=\sqrt2\)，\(\Longrightarrow\)，\(a-\frac14=2\)，\(\Longrightarrow\)，\(a=\frac94\).
\end{solution}

\section{解答题}

\begin{problem}
在 \(\triangle ABC\) 中，内角 \(A\)，\(B\)，\(C\) 的对边分别为 \(a\)，\(b\)，\(c\)，且 \(b \sin A = \sqrt{3} a \cos B\).

\begin{enumerate}
\item 求角 \(B\) 的大小；
\item 若 \(b = 3\)，\( \sin C = 2 \sin A \)，求 \(a\)，\(c\) 的值.
\end{enumerate}
\end{problem}

\begin{solution}
\begin{enumerate}
\item 由正弦定理 \(b/a=\sin B/\sin A\)，题设条件化为
\[
a\sin B=\sqrt3a\cos B
\]
因为 \(a>0\)，所以
\[
\sin B=\sqrt3\cos B
\]
三角形内角 \(B\in(0,\pi)\)，上式右端必须为正，故 \(\cos B>0\)，从而
由此 \(\tan B=\sqrt3\)，所以 \(B=\frac{\pi}{3}\).

\item 再由正弦定理和 \(\sin C=2\sin A\)，有
\[
\frac ca=\frac{\sin C}{\sin A}=2
\]
所以 \(c=2a\). 由余弦定理及 \(\cos B=1/2\)，
\[
b^2=a^2+c^2-2ac\cos B
=a^2+4a^2-2a\cdot2a\cdot\frac12=3a^2
\]
因 \(b=3\)，故 \(a=\sqrt3\)，\(c=2\sqrt3\).
\end{enumerate}
\end{solution}

\begin{problem}
已知数列 \(\{a_{n}\}\) 的前 \(n\) 项和为 \(S_{n}\)，且 \(S_{n} = 2n^{2} + n\)，\(n \in \N^{*}\)，数列 \(\{b_{n}\}\) 满足 \(a_{n} = 4\log_{2}b_{n} + 3\)，\(n \in \N^{*}\).

\begin{enumerate}
\item 求 \(a_{n}\)，\(b_{n}\)；
\item 求数列 \(\{a_{n} \cdot b_{n}\}\) 的前 \(n\) 项和 \(T_{n}\).
\end{enumerate}
\end{problem}

\begin{solution}
\begin{enumerate}
\item 先由 \(S_1=3\) 得 \(a_1=3=4\cdot1-1\). 当 \(n\geqslant2\) 时，
\[
a_n=S_n-S_{n-1}
=(2n^2+n)-\bigl(2(n-1)^2+(n-1)\bigr)=4n-1
\]
因此对一切 \(n\in\N^*\)，
\[
a_n=4n-1
\]
代入 \(a_n=4\log_2b_n+3\)，得
\(\log_2b_n=n-1\)，所以 \(b_n=2^{n-1}\).

\item 于是
\[
T_n=\sum_{k=1}^n(4k-1)2^{k-1}
\]
令 \(u_k=(4k-5)2^k\)，则
\[
u_k-u_{k-1}
=(4k-5)2^k-(4k-9)2^{k-1}
=(4k-1)2^{k-1}
\]
所以
\[
T_n=u_n-u_0=(4n-5)2^n-(-5)=(4n-5)2^n+5
\]
\end{enumerate}
\end{solution}

\begin{problem}
如图，在侧棱垂直底面的四棱柱 \(ABCD - A_{1}B_{1}C_{1}D_{1}\) 中，\(AD \parallel BC\)，\(AD \perp AB\)，\(AB = \sqrt{2}\)，\(AD = 2\)，\(BC = 4\)，\(AA_{1} = 2\)，\(E\) 是 \(DD_{1}\) 的中点，\(F\) 是平面 \(B_{1}C_{1}E\) 与直线 \(AA_{1}\) 的交点.

\begin{enumerate}
    \item 证明：
    \begin{enumerate}
    \item \(EF \parallel A_{1}D_{1}\)；
\item \(BA_{1} \perp\) 平面 \(B_{1}C_{1}EF\).
\end{enumerate}
\item 求 \(BC_{1}\) 与平面 \(B_{1}C_{1}EF\) 所成的角的正弦值.
\end{enumerate}

\bitmapfigure[max width=0.6\linewidth]{img/2012/zhejiang_liberal/q20_fig1.png}
\end{problem}

\begin{solution}
建立空间直角坐标系，取 \(A=(0,0,0)\)，\(B=(\sqrt2,0,0)\)，\(D=(0,2,0)\)，\(C=(\sqrt2,4,0)\).
由侧棱垂直底面且 \(AA_1=2\)，得
\(A_1=(0,0,2)\)，\(B_1=(\sqrt2,0,2)\)，\(C_1=(\sqrt2,4,2)\)，\(D_1=(0,2,2)\).
又 \(E\) 是 \(DD_1\) 的中点，故 \(E=(0,2,1)\).

\begin{enumerate}
\item 证明.
\begin{enumerate}
\item 平面 \(B_1C_1E\) 的一个方程为
\[
x-\sqrt2z+\sqrt2=0
\]
因为 \(B_1\)，\(C_1\)，\(E\) 的坐标都满足此方程. 直线 \(AA_1\) 上 \(x=0\)，代入得 \(z=1\)，所以
\[
F=(0,0,1)
\]
于是
\(\overrightarrow{EF}=(0,-2,0)\)，\(\overrightarrow{A_1D_1}=(0,2,0)\)，
故 \(EF\parallel A_1D_1\).

\item 再有
\(\overrightarrow{BA_1}=(-\sqrt2,0,2)\)，\(\overrightarrow{B_1C_1}=(0,4,0)\)，\(\overrightarrow{EF}=(0,-2,0)\).
向量 \(\overrightarrow{BA_1}\) 垂直于 \(\overrightarrow{B_1C_1}\)，且
\[
\overrightarrow{BA_1}\cdot\overrightarrow{B_1F}
=(-\sqrt2,0,2)\cdot(-\sqrt2,0,-1)=0
\]
直线 \(B_1C_1\) 与 \(B_1F\) 都在平面 \(B_1C_1EF\) 内，且相交于 \(B_1\)，故
\[
BA_1\perp\text{平面 }B_1C_1EF
\]

\end{enumerate}
\item 由于 \(BC_1\) 的方向向量为
\[
\overrightarrow{BC_1}=(0,4,2)
\]
且 \(\overrightarrow{BA_1}\) 是该平面的法向量，直线与平面所成角的正弦为
\[
\frac{|\overrightarrow{BC_1}\cdot\overrightarrow{BA_1}|}
{|\overrightarrow{BC_1}|\,|\overrightarrow{BA_1}|}
=\frac4{\sqrt{20}\sqrt6}
=\frac{\sqrt{30}}{15}
\]
\end{enumerate}
\end{solution}

\begin{problem}
已知 \(a \in \R\)，函数 \(f(x) = 4x^3 - 2ax + a\).

\begin{enumerate}
\item 求 \(f(x)\) 的单调区间；
\item 证明：当 \(0 \leqslant x \leqslant 1\) 时，\(f(x) + |2 - a| > 0\).
\end{enumerate}
\end{problem}

\begin{solution}
\begin{enumerate}
\item
\[
f'(x)=12x^2-2a
\]
当 \(a\leqslant0\) 时，\(f'(x)\geqslant0\)，且除至多一个点外为正，所以 \(f(x)\) 在 \(\R\) 上单调增加.

当 \(a>0\) 时，令
\[
r=\sqrt{\frac a6}
\]
则 \(f'(x)>0\) 当 \(x<-r\) 或 \(x>r\)，\(f'(x)<0\) 当 \(-r<x<r\). 因此增区间为 \(\left(-\infty,-\sqrt{\frac a6}\right]\) 和 \(\left[\sqrt{\frac a6},+\infty\right)\)，减区间为 \(\left[-\sqrt{\frac a6},\sqrt{\frac a6}\right]\).

\item 令
\[
g(x)=2x^3-2x+1
\]
在 \([0,1]\) 上，
\[
g'(x)=6x^2-2
\]
在 \([0,1/\sqrt3)\) 上为负，在 \((1/\sqrt3,1]\) 上为正，因此其最小值在 \(x=1/\sqrt3\) 处取得，且
\[
g\left(\frac1{\sqrt3}\right)
=1-\frac4{3\sqrt3}>0
\]
又因为 \(3\sqrt3>4\)，所以 \(1-\frac{4}{3\sqrt3}>0\)，故 \(g(x)>0\).

若 \(a\leqslant2\)，则
\[
f(x)+|2-a|=4x^3-2ax+2
\geqslant4x^3-4x+2=2g(x)>0
\]
其中用到了 \(x\geqslant0\) 及 \(a\leqslant2\). 若 \(a>2\)，则
\[
f(x)+|2-a|=4x^3+2a(1-x)-2
\geqslant4x^3+4(1-x)-2
=4x^3-4x+2=2g(x)>0
\]
其中用到了 \(1-x\geqslant0\). 两种情形均得证.
\end{enumerate}
\end{solution}

\begin{problem}
如图，在直角坐标系 \(xOy\) 中，点 \(P\left(1, \frac{1}{2}\right)\) 到抛物线 \(C: y^2 = 2px \) （\(p > 0\)）的准线的距离为 \(\frac{5}{4}\). 点 \(M(t, 1)\) 是 \(C\) 上的定点，\(A\)，\(B\) 是 \(C\) 上的两动点，且线段 \(AB\) 被直线 \(OM\) 平分.

\begin{enumerate}
\item 求 \(p\)，\(t\) 的值；
\item 求 \(\triangle ABP\) 面积的最大值.
\end{enumerate}

\bitmapfigure[max width=0.28\linewidth]{img/2012/zhejiang_liberal/q22_fig1.png}
\end{problem}

\begin{solution}
抛物线 \(y^2=2px\) 的准线为 \(x=-p/2\). 点 \(P\) 到准线的距离为
\[
1+\frac p2=\frac54
\]
所以 \(p=\frac12\)，抛物线方程为 \(y^2=x\). 由于 \(M(t,1)\) 在抛物线上，故
\[
t=1
\]
此时直线 \(OM\) 的方程为 \(y=x\).

设 \(A=(y_1^2,y_1)\)，\(B=(y_2^2,y_2)\)，
并令 \(s=y_1+y_2\)，\(\ q=y_1y_2\). 线段 \(AB\) 被 \(y=x\) 平分，意味着其中点在 \(y=x\) 上，故
\[
\frac{y_1^2+y_2^2}{2}=\frac{y_1+y_2}{2}
\]
即
\(y_1^2+y_2^2=s\)，所以 \(q=\frac{s^2-s}{2}\).
由于 \(A\ne B\)，二次方程 \(y^2-sy+q=0\) 有两个不等实根，故
\[
s^2-4q=s(2-s)>0
\]
所以 \(0<s<2\).

弦 \(AB\) 的方程为
\[
x-sy+q=0
\]
又
\[
|AB|=\sqrt{s(2-s)}\sqrt{1+s^2}
\]
点 \(P(1,1/2)\) 到弦所在直线的距离为
\[
\frac{|1-\frac s2+q|}{\sqrt{1+s^2}}
=\frac{(s-1)^2+1}{2\sqrt{1+s^2}}
\]
其中用到了 \(q=(s^2-s)/2\). 因此三角形面积为
\[
S=\frac12|AB|\cdot d(P,AB)
=\frac14\sqrt{s(2-s)}\bigl((s-1)^2+1\bigr)
\]
令 \(u=s-1\)，则 \(|u|<1\)，从而
\[
S=\frac14(1+u^2)\sqrt{1-u^2}
\]
再令 \(v=u^2\in[0,1)\)，则
\[
S^2=\frac1{16}(1+v)^2(1-v)
\]
函数 \(h(v)=(1+v)^2(1-v)\) 的导数为
\[
h'(v)=(1+v)(1-3v)
\]
所以 \(h\) 在 \(0\leqslant v\leqslant1/3\) 上增加，在 \(1/3\leqslant v<1\) 上减少，最大值在 \(v=1/3\) 处取得. 于是
\[
S_{\max}
=\frac14\left(1+\frac13\right)\sqrt{1-\frac13}
=\frac{\sqrt6}{9}
\]
\end{solution}
